\documentclass[12pt,a4paper]{article}

\usepackage{kotex}
\usepackage{amsmath,amssymb,amsfonts}
\usepackage{graphicx}
\usepackage{textcomp}
\usepackage{xcolor}
\usepackage{booktabs}
\usepackage{multirow}
\usepackage{url}
\usepackage{array}
\usepackage{balance}
\graphicspath{{../../figures/}}

\begin{document}

\title{TA-BRPL 연구노트 v0.3.3\\
Blackhole + Sinkhole 복합 공격 환경에서의\\
Trust-Aware Backpressure RPL 구현 및 초기 결과}

\author{연구노트 v0.3.3 --- Internal Draft\\최종 수정: 2026-03-19}

\maketitle

% ============================================================
\begin{abstract}
% ============================================================
본 문서는 TA-BRPL(Trust-Aware Backpressure RPL)의 현재 구현 상태와 최신 실험 결과를
통합 정리한 연구노트이다. 대상 환경은 Contiki-NG/Cooja 기반의 $6\times6$ GRID 토폴로지
(총 36개 노드, 30개 random seed)이며, 3개의 blackhole 공격자(node 2, 3, 4)와
1개의 sinkhole 공격자(node 18)가 동시에 동작하는 복합 공격 시나리오를 가정한다.

TA-BRPL은 BRPL의 backpressure 기반 경로 선택 능력을 유지하면서,
3성분 on-node trust 모델($T_\text{fwd}$, $T_\text{ctrl}$, $T_\text{hon}$)을
BRPL objective function에 직접 결합한다. 전체 30-seed 실험에서 TA-BRPL은
during-attack PDR \textbf{0.9077}, recovery PDR \textbf{0.9593}을 기록하여,
RPL(0.8759/0.8754), BRPL(0.8508/0.8399), SMTrust(0.8701/0.8649) 대비
더 높은 공격 복원력을 보였다.

v0.3.3에서는 기존 30-seed 기본 결과를 유지한 상태에서, 링크 손실 강건성 분석을 확장하였다.
구체적으로 ETX 기반 PRR 정규화, 절대 threshold 재튜닝($\tau_\text{join}=580$, $\tau_\text{warn}=630$),
soft-penalty sweep, 그리고 median 기반 relative trust filtering을 구현하고 5-seed loss ladder
(10\%, 30\%, 50\%)에서 비교하였다. 그 결과, 강한 soft penalty는 공격자 사용 비율은 줄이지만
정상 parent까지 과도하게 pruning하는 문제가 확인되었고, relative filtering은 absolute threshold보다
``attacker share 감소 + usable parent 유지''의 균형을 더 잘 달성하였다.

본 문서는 설계 의도, 전체 구현 구조, 실험 설정, 30-seed 기본 결과,
EWMA/임계값 민감도 분석, 성분 절제 실험, 혼잡--공격 분리 실험,
오탐 실험, 링크 손실 robustness와 tuning progression, 그리고 현재 한계와 후속 과제를 정리한다.
\end{abstract}

\textbf{키워드:} IoT, LLN, RPL, BRPL, Trust Management, Blackhole, Sinkhole,
Selective Forwarding, Cooja, Contiki-NG

% ============================================================
\section{연구 목표}
% ============================================================

본 연구의 목적은 LLN(Low-power and Lossy Network) 환경에서 다음의 두 가지 상충되는 요구를
동시에 만족시키는 것이다. 첫째, BRPL이 제공하는 backpressure 기반 혼잡 인식 라우팅
능력을 유지한다. 둘째, blackhole(선택적 패킷 드롭)과 sinkhole(rank 조작)과 같은
insider routing attack에 대해 노드 내 자율적 복원력을 확보한다.

기존 접근의 한계는 크게 두 가지다. 첫째, 단순 forwarding rate 기반 감시는
혼잡과 공격을 구분하지 못한다. 즉, 정직하지만 혼잡한 노드와 의도적으로 패킷을
드롭하는 공격자는 모두 forwarding rate 하락을 유발한다. 둘째, SMTrust와 같이
weighted-sum trust를 MRHOF 위에 얹는 구조는 정적 환경에서 일부 가중치 성분이
상수처럼 고정되어 trust floor가 탐지 임계값 이상에 묶이는 구조적 한계를 가진다
(Section~\ref{sec:smtrust} 참조).

TA-BRPL의 핵심 가설은 다음과 같다.
\begin{itemize}
  \item $T_\text{fwd}$: 수동 도청(passive overhearing) 기반 forwarding success rate로 blackhole을 탐지한다.
  \item $T_\text{ctrl}$: DIO rank/version 이상 징후를 활용해 sinkhole을 탐지한다.
  \item $T_\text{hon}$: 큐 광고값과 로컬 추정치의 교차 검증을 통해 혼잡과 공격을 분리한다.
  \item 세 성분의 가중 기하평균을 trust로 집계하고 이를 BRPL cost function에 직접 결합함으로써,
  공격자 parent가 점진적으로 배제되도록 한다.
\end{itemize}

% ============================================================
\section{전체 시스템 구조}
% ============================================================

\subsection{계층 구조}

TA-BRPL은 Contiki-NG의 RPL-Classic 라우팅 스택 위에서 동작한다.
핵심 모듈은 \texttt{ta-brpl-trust.c/h}(trust 엔진, 약 850 LoC)이며,
\texttt{rpl-brpl.c}의 objective function에 C weak-symbol 오버라이드를 통해
결합된다. 이 설계는 TA-BRPL을 컴파일에서 제거하더라도 BRPL/RPL의
원본 소스를 직접 수정할 필요가 없음을 의미한다.

네트워크 스택에서 trust 엔진으로 입력되는 훅은 다음 세 가지다.
\begin{enumerate}
  \item \textbf{IP input hook} (\texttt{NETSTACK\_CONF\_IP\_INPUT\_CALLBACK}):  
  릴레이 unicast 패킷을 감지하여 $T_\text{fwd}$ 관찰 이벤트를 공급하고,
  DIO ICMPv6 메시지를 파싱하여 $T_\text{ctrl}$ 이상 지표를 갱신한다.
  \item \textbf{BRPL queue callback}:  
  이웃 노드의 큐 광고값($Q_\text{adv}$, $Q_\text{max}$)을 전달하여
  $T_\text{hon}$ 계산에 사용한다.
  \item \textbf{RPL parent-change callback}
  (\texttt{RPL\_CALLBACK\_PARENT\_SWITCH}):  
  preferred parent 교체 시 \texttt{current\_parent\_id}와 escape 타이머를 갱신한다.
  이 훅이 없으면 escape 메커니즘은 발동하지 않는다.
\end{enumerate}

trust 엔진에서 BRPL 측으로 전달되는 출력은 다음의 세 weak-symbol 함수로 제공된다.
\begin{itemize}
  \item \texttt{brpl\_trust\_get(node\_id)}: trust 값 [0, 1000] 반환
  \item \texttt{brpl\_penalty\_scale\_get(node\_id)}: 지속 패널티 배율 반환
  \item \texttt{brpl\_escape\_mode\_get(node\_id)}: 탈출 모드 플래그 반환
\end{itemize}

\subsection{파라미터 단일 소스}

모든 tunable parameter는 \texttt{project-conf.h}의 C 매크로로 정의된다.
각 프로토콜 변형의 \texttt{Makefile.*}은 모드 플래그만 설정하며,
임계값 및 $\lambda$ 값은 모두 \texttt{project-conf.h}에서만 관리된다.
민감도 실험은 컴파일러 커맨드라인에서 특정 매크로만 오버라이드하는 방식으로 수행한다.

% ============================================================
\section{Trust 모델 상세}
% ============================================================

\subsection{$T_\text{fwd}$: Forwarding Trust}

$T_\text{fwd}$는 이웃 $j$가 패킷을 얼마나 신뢰성 있게 중계하는지를
측정하는 성분이다. 노드 $i$는 이웃 $j$를 통해 패킷을 전송할 때마다
카운터 $S_{ij}$를 증가시키고, root로부터 echo reply가 돌아오면
해당 TX를 실제로 실어 나른 parent $j$에 대해 forwarding 관측치 $F_{ij}$를 증가시킨다.
초기 trust가 과도하게 0 또는 1로 치우치지 않도록 Laplace smoothing
파라미터 $(\alpha,\beta)$를 적용하여 초기값을 0.5 근방으로 유지한다.
\begin{equation}
  T_\text{fwd}(i,j) = \frac{F_{ij} + \alpha}{E_{ij} + \alpha + \beta}, \quad
  E_{ij} = S_{ij}\,\widehat{\mathrm{PRR}}_{ij}.
  \label{eq:tfwd}
\end{equation}

여기서 $\widehat{\mathrm{PRR}}_{ij}$는 Contiki-NG의 \texttt{link-stats} 모듈이 유지하는
이웃별 ETX로부터 추정한다.
\begin{equation}
  \widehat{\mathrm{PRR}}_{ij}
  = \mathrm{clamp}\!\left(
      \frac{\texttt{ETX\_DIVISOR}}{\texttt{etx}_{ij}},
      \mathrm{PRR}_{\min},
      1
    \right).
\end{equation}
ETX가 아직 수렴하지 않았거나 freshness가 부족한 경우에는
$\widehat{\mathrm{PRR}}_{ij}=1$로 두어 join 직후 FP를 방지한다.

카운터는 60초 업데이트 주기마다 \textbf{우측 비트시프트(halving)}로 감쇠된다.
여기서 리셋 대신 halving을 사용하는 이유는, 완전 리셋을 사용할 경우
공격자가 잠시 정상적으로 패킷을 전달하는 것만으로 trust가 다시 0.5 수준으로
중립화되는 문제가 발생하기 때문이다.

\subsection{$T_\text{ctrl}$: Control-Plane Trust}

$T_\text{ctrl}$은 제어 평면 이상을 세 가지 지표로 탐지한다.

\textbf{Rank 이상 $A_\text{rank}$}: DIO에 포함된 rank가 root rank보다 낮으면
물리적으로 불가능한 상태이므로 즉시 최대 이상으로 처리한다.
또한 과도한 rank 변동은 hysteresis 기준을 초과한 변화 횟수로 누적한다.

\textbf{DIO 빈도 이상 $A_\text{DIO}$}: trickle 타이머 기준의 기대 전송률 대비
과도하게 많은 DIO 전송 횟수를 누적하여 DIO suppression/flooding 이상을 탐지한다.

\textbf{Version 이상 $A_\text{ver}$}: root에서 전파되지 않은 RPL version 번호의 증가는
DODAG 강제 리셋을 유발하는 version attack의 신호로 해석한다.

세 지표를 가중합하여 제어 평면 이상 점수를 계산하고 이를 trust로 변환한다.
\begin{equation}
  A_{ij} = \frac{w_1 A_\text{rank} + w_2 A_\text{DIO} + w_3 A_\text{ver}}{10},\quad
  T_\text{ctrl}(i,j) = 1 - A_{ij}.
\end{equation}
기본 가중치는 $w_1=5$, $w_2=3$, $w_3=2$이다.

\subsection{$T_\text{hon}$: Congestion-Honesty Trust}

$T_\text{hon}$은 이웃 $j$가 자신의 큐 상태를 정직하게 광고하는지를
검증하는 성분이다.
\begin{equation}
  T_\text{hon}(i,j) = 1 - \min\!\left(1,\; \frac{|Q_\text{adv} - Q_\text{est}|}{Q_\text{max}}\right).
\end{equation}

여기서 $Q_\text{adv}$는 이웃 $j$가 BRPL 비콘에 실어 보낸 큐 점유율이고,
$Q_\text{est}$는 노드 $i$가 자신의 로컬 상태로부터 추정한
이웃 $j$의 예상 부하다. 정직하게 혼잡한 노드는 대체로
$Q_\text{adv} \approx Q_\text{est}$를 만족하므로 $T_\text{hon} \approx 1.0$을 유지한다.
반면 공격자가 큐 상태를 실제보다 낮게 광고하면 두 값의 차이가 커지고,
그 결과 $T_\text{hon}$이 하락한다.

이 성분은 ``혼잡으로 인해 forwarding rate가 낮아진 정직한 노드''와
``의도적으로 패킷을 드롭하는 공격자''를 구분하기 위해 필요하다.

\subsection{가중 기하평균 집계}

세 trust 성분은 다음과 같이 가중 기하평균으로 결합된다.
\begin{equation}
  \tilde{T}(i,j) =
  T_\text{fwd}^{w_f/W}
  \cdot
  T_\text{ctrl}^{w_c/W}
  \cdot
  T_\text{hon}^{w_h/W},
\end{equation}
여기서 $W = w_f + w_c + w_h = 10$이며, 기본 가중치 벡터는
$(w_f, w_c, w_h) = (7, 2, 1)$이다.
현재 구현은 forwarding 신호에 더 큰 비중을 두어,
공격자 배제 이후에도 $T_\text{fwd}$ 하락이 aggregate trust에 충분히 반영되도록 설계되어 있다.

기하평균을 사용하는 이유는 어느 한 성분이라도 0에 가까워지면
전체 trust가 크게 하락하기 때문에, 공격자가 특정 단일 지표만 조작해서
높은 trust를 유지하기 어렵기 때문이다.

\subsection{비대칭 EWMA 업데이트}

60초마다 다음 점화식으로 per-neighbor trust를 업데이트한다.
\begin{equation}
  T_{ij}(t+1) = \lambda(t)\, T_{ij}(t) + (1-\lambda(t))\, \tilde{T}_{ij},
\end{equation}
\begin{equation}
  \lambda(t) =
  \begin{cases}
    \lambda_\downarrow = 0.2, & \tilde{T}_{ij} < T_{ij}(t)\quad \text{(신뢰 하락 시)} \\
    \lambda_\uparrow   = 0.7, & \text{otherwise}\quad \text{(신뢰 유지/회복 시).}
  \end{cases}
\end{equation}

$\lambda_\downarrow=0.2$는 공격 직후 trust가 빠르게 하락하도록 허용하고,
$\lambda_\uparrow=0.7$은 회복 속도를 의도적으로 늦춰,
공격자가 잠시 정상 동작을 보이며 trust를 회복하는 on-off 전략에 대응한다.

모든 연산은 고정소수점(정수 범위 $[0,1000]$, 해상도 1/1000)으로 수행된다.
기하평균 계산에 필요한 지수 연산은 추후 lookup table 기반으로 대체 가능하다.

\subsection{임계값 정책과 튜닝 경로}

\begin{table}[t]
  \centering
  \caption{TA-BRPL trust tier 정의(기본선)}
  \label{tab:thresholds}
  \begin{tabular}{lllp{3.2cm}}
    \toprule
    Tier & 범위 & 상태 & BRPL 효과 \\
    \midrule
    Normal    & $T \ge 700$       & NORMAL      & 패널티 없음 \\
    Suspect   & $450 \le T < 700$ & SUSPECT     & cost 배율 패널티 적용 \\
    Untrusted & $250 \le T < 450$ & UNTRUSTED   & parent 후보 제외 \\
    Blacklist & $T < 250$         & BLACKLISTED & 하드 격리 \\
    \bottomrule
  \end{tabular}
\end{table}

parent 선택 시 이웃 $j$의 BRPL routing cost는 다음과 같이 조정된다.
\begin{equation}
  \text{cost}'(j) = \text{cost}(j) \times \frac{P_s(j)}{1000},
\end{equation}
여기서 $P_s(j)$는 Normal 노드에서는 1000(즉, 효과 없음)이며,
Suspect 노드에서는 상황에 따라 가변 패널티를 갖는다.
본선 30-seed 기본 결과는 이 절대 임계값 구조를 사용하였다.
반면 손실 robustness tuning에서는
$(\tau_\text{join},\tau_\text{warn})=(580,630)$과 같은 재튜닝 및
relative trust filtering 변형을 추가로 탐색하였다.

\subsection{지속 패널티와 탈출 메커니즘}

현재 preferred parent가 공격자인 경우, RPL 자체의 parent-change hysteresis가
전환을 지연시키거나 막을 수 있다. TA-BRPL은 이를 극복하기 위해 두 가지 메커니즘을 둔다.

\textbf{Attack persistence penalty}: preferred parent의 trust가
$\tau_\text{warn}$ 미만인 상태가 지속되면, 매 120초 창마다
패널티 배율에 250을 가산한다(최대 2600).

\textbf{Escape mode}: preferred parent의 trust가 $\tau_\text{warn}$ 미만인 상태가
180초 이상 지속되면 escape mode를 활성화한다.
이 모드에서는 현재 parent에 대한 hysteresis가 해제되고,
추가 패널티 배율 3200이 적용되어 metric 차이와 무관하게 parent 교체가 강제된다.
현재 코드에서는 여기에 더해
\texttt{low-trust} 연속 업데이트 조건,
$T_\text{fwd}<450$ 및 $T_\text{hon}>850$의 원인 기반 조건,
그리고 선택적으로 better-parent-only 조건을 결합할 수 있도록 구현되어 있다.

격리 해제 후에는 trust를 $\tau_\text{join}$으로 복원하고,
패널티 배율을 120초 동안 $1.6\times \rightarrow 1.0\times$로 선형 감쇠시키며,
재드롭 감시 플래그를 활성화하여 공격이 재개될 경우 즉시 재격리한다.
또한 현재 구현은 join/blacklist 판정에 최소 지속 시간(기본 120초)을 적용하여,
일시적 혼잡이나 단기 loss dip이 즉시 hard exclusion으로 이어지지 않도록 한다.

v0.3.3의 tuning 실험에서는 escape churn을 줄이기 위해
연속 low-trust 조건, cooldown, better-parent-only 조건을 포함한 Phase 3 변형도 탐색하였다.
이 변형들은 attacker share와 parent churn을 줄이는 데는 유효했지만,
단독으로는 BRPL과의 PDR 격차를 해소하지 못했다.

\subsection{Loss-Robustness 확장: Relative Trust Filtering}

v0.3.3에서는 손실 환경에서 절대 임계값만으로 parent 후보를 자를 경우
정상 이웃까지 함께 제외되는 문제를 완화하기 위해,
\textbf{median 기반 relative filtering} 실험 변형을 추가하였다.
핵심 아이디어는 ``$T < \tau_\text{join}$이면 제외''라는 절대 기준을
현재 이웃 trust 분포에 대한 상대 기준으로 일부 대체하는 것이다.

\begin{equation}
  T^\text{floor}_i = \max\!\bigl(\tau_\text{black},\; \mathrm{median}_i(T) - \Delta_\text{rel}\bigr),
\end{equation}
\begin{equation}
  j \in \mathcal{P}_i \iff T_{ij} \ge T^\text{floor}_i.
\end{equation}

여기서 $\Delta_\text{rel}$은 relative margin이다. 이 변형에서는 공격자를
``절대적으로 낮은 trust''가 아니라, \emph{이웃 집단 대비 유의미하게 낮은 trust}를
보이는 노드로 해석한다. 따라서 링크 손실로 인해 모든 이웃의 trust가 함께 하락해도,
정상 노드들이 median 근방에 모여 있으면 후보 집합이 과도하게 축소되지 않는다.

soft penalty 역시 절대 trust 값이 아니라 median 대비 gap에 따라 부여하였다.
\begin{equation}
  \mathrm{gap}_{ij} = \max(0,\; \mathrm{median}_i(T) - T_{ij}),
\end{equation}
\begin{equation}
  P^\text{rel}_{ij} = \min\!\left(P_\text{max},\;
  \frac{\mathrm{gap}_{ij}}{1000}\, S_\text{rel}\right).
\end{equation}

이 실험 변형은 본선 30-seed 기본 결과를 대체하는 최종 기본 설정은 아니며,
loss ladder에서 absolute threshold 구조의 한계를 진단하고
후속 설계를 탐색하기 위한 \emph{robustness-oriented tuned variant}로 사용하였다.

% ============================================================
\section{SMTrust 구현 상태 및 구조적 한계}
\label{sec:smtrust}
% ============================================================

\subsection{현재 구현}

SMTrust는 ``RPL + trust-aware parent filtering'' 구조로 구현되어 있다.
기반 라우팅은 MRHOF이며, parent 선택 시 다음의 두 단계를 거친다.
\begin{enumerate}
  \item trust class $t_1$/$t_2$ 노드를 후보에서 제외한다 (threshold 0.46).
  \item near-tie 상황에서는 $t_4$/$t_5$ 노드를 $t_3$ 노드보다 우선한다.
\end{enumerate}

Trust index는 다음과 같은 6개 성분의 가중합으로 계산된다.
\begin{equation}
  \text{TI} = 0.25\,\text{TMSR} + 0.15\,\text{TM}_{H_0}
             + 0.15\,\text{TMEL} + 0.20\,\text{TMLLS}
             + 0.15\,\text{TMMobility} + 0.10\,\text{TMRT}.
\end{equation}

\subsection{구조적 한계: TI Floor}

현재 실험 환경은 정적 $6\times6$ GRID 토폴로지이므로, 두 성분이 사실상 상수로 고정된다.
\begin{itemize}
  \item \textbf{TMEL} (가중치 0.15): 이웃의 실제 잔여 에너지가 아니라
  로컬 Energest 추정치를 사용하므로 환경 전반에서 거의 상수 $\approx 0.50$로 유지된다.
  \item \textbf{TMMobility} (가중치 0.15): 정적 토폴로지이므로 $=1.00$으로 고정된다.
\end{itemize}

즉, 총 결합 가중치 0.30이 변별력을 갖지 못한 채 상수처럼 작용한다.
$\text{TMSR}=0$인 완전 blackhole 극단 시나리오를 가정하면 TI의 최솟값은 다음과 같다.
\begin{align}
  \text{TI}_\text{min}
  &= 0 \cdot 0.25 + 0.50 \cdot 0.15 + 0.50 \cdot 0.15 \notag\\
  &\quad + 0.57 \cdot 0.20 + 1.00 \cdot 0.15 + 0.50 \cdot 0.10 \notag\\
  &= 0.464 > 0.460 \quad \text{(필터링 임계값)}.
\end{align}
30-seed 실측 최솟값 역시 $\text{TI}_\text{min}=0.489$였다.

따라서 공격자 노드는 아무리 패킷을 드롭하더라도 TI가 임계값 아래로
충분히 내려가지 않아 parent 후보 집합에서 제거되지 않는다.

이는 단순 구현 버그라기보다, \textbf{구현--환경 불일치}
(implementation-environment mismatch)로 해석하는 것이 더 적절하다.
즉, 실제 이동성이 존재하거나 에너지 정보 교환이 가능한 환경에서는
TMMobility와 TMEL이 변별력을 가질 수 있으나, 현재 실험 환경에서는 그렇지 않다.

% ============================================================
\section{실험 설정}
% ============================================================

\subsection{토폴로지와 노드 역할}

실험 토폴로지는 $6\times6$ GRID, 총 36개 노드이며,
노드 간 간격은 40\,m, 전파 모델은 UDGM
(\texttt{tx\_range=50\,m}, \texttt{success\_ratio=1.0})를 사용하였다.

\begin{itemize}
  \item \textbf{Node 1}: DODAG root (border router)
  \item \textbf{Node 2, 3, 4}: Blackhole 공격자. 350초 이후 전달해야 할 UDP 패킷을 100\% 드롭하며, DIO/routing 동작은 정상적으로 유지
  \item \textbf{Node 18}: Sinkhole 공격자. DIO에 rank$+1$을 광고하여 인근 노드를 유인하고, 유입된 트래픽은 드롭
  \item \textbf{나머지 31개 노드}: Sender. 30초마다 root에 UDP 패킷을 전송
\end{itemize}

\begin{figure}[t]
  \centering
  \includegraphics[width=0.95\linewidth]{fig0_topology_grid6x6.pdf}
  \caption{실험 토폴로지($6\times6$ GRID). Node 1은 root, Node 2/3/4는 blackhole,
  Node 18은 sinkhole, 나머지 노드는 sender이다.}
  \label{fig:topology_grid}
\end{figure}

\subsection{실험 타임라인}

\begin{itemize}
  \item $[0,150)$초: Warm-up. DODAG 형성 및 초기 수렴 단계, 통계 미수집
  \item $[150,350)$초: Pre-attack. 정상 동작 구간, 기준 PDR 측정
  \item $[200,300)$초: \textit{(pre-attack 내부)} 합성 혼잡 버스트.
  root 인접 5개 노드가 15초 주기로 추가 송신하여 큐 포화를 유도하고,
  $T_\text{hon}$의 혼잡 인식 능력을 검증
  \item $[350,650)$초: During attack. 공격자 활성화, 주요 PDR 측정 구간
  \item $[650,900)$초: Recovery. 공격자는 계속 활성 상태를 유지하나,
  trust 기반 우회가 수렴한 이후의 PDR 측정
\end{itemize}

전체 시뮬레이션 시간은 900초이며, random seed는 총 30개를 사용하였다.
통계 검정은 Wilcoxon rank-sum test(비모수, Bonferroni 보정)를 적용하였다.

\subsection{비교 프로토콜}

\begin{itemize}
  \item \textbf{RPL}: RPL-Classic + MRHOF/ETX, trust 비활성
  \item \textbf{BRPL}: BRPL + queue-aware cost, trust 비활성
  \item \textbf{SMTrust}: RPL-Classic MRHOF + 6성분 trust filtering
  \item \textbf{TA-BRPL}: 제안 프로토콜
\end{itemize}

% ============================================================
\section{전체 30-Seed 기본 결과}
\label{sec:main}
% ============================================================

Table~\ref{tab:main_pdr}는 4개 프로토콜의 phase별 평균 PDR을 보여준다.

\begin{table}[t]
  \centering
  \caption{Phase별 평균 PDR (30 seeds)}
  \label{tab:main_pdr}
  \begin{tabular}{lrrr}
    \toprule
    프로토콜 & Pre-attack & During attack & Recovery \\
    \midrule
    TA-BRPL  & 0.9993 & \textbf{0.9077} & \textbf{0.9593} \\
    RPL      & 1.0000 & 0.8759 & 0.8754 \\
    SMTrust  & 0.9998 & 0.8701 & 0.8649 \\
    BRPL     & 0.9999 & 0.8508 & 0.8399 \\
    \bottomrule
  \end{tabular}
\end{table}

\textbf{TA-BRPL vs. RPL}: during에서 $+3.2$\,pp, recovery에서 $+8.4$\,pp 개선되었다.  
\textbf{TA-BRPL vs. BRPL}: during에서 $+5.7$\,pp, recovery에서 $+11.9$\,pp 개선되었다.  
모든 pairwise 비교에서 Wilcoxon 검정 결과는 $p < 0.001$이었다.

\begin{figure}[t]
  \centering
  \includegraphics[width=\linewidth]{fig1_pdr_phases.pdf}
  \caption{Phase별 PDR 비교 (30 seeds). 각 박스는 25--75 백분위수, 수염은 10--90 백분위수를 나타낸다.
  TA-BRPL은 during-attack과 recovery 두 구간 모두에서 가장 높은 PDR을 달성한다.}
  \label{fig:pdr_phases}
\end{figure}

\begin{figure}[t]
  \centering
  \includegraphics[width=\linewidth]{fig2_resilience_summary.pdf}
  \caption{공격 복원력 요약. TA-BRPL에서는 recovery PDR이 during-attack PDR보다 높게 나타나는
  trust convergence effect가 뚜렷하게 관찰된다.}
  \label{fig:resilience}
\end{figure}

\subsection{Trust Convergence Effect}

TA-BRPL의 recovery PDR(0.9593)은 during-attack PDR(0.9077)보다 높다.
공격자가 recovery 구간에서도 계속 활성 상태임을 고려하면 이는 직관에 반하는 결과처럼 보인다.

그 원인은 150초 주기의 EWMA 업데이트 지연에 있다.
\begin{itemize}
  \item 1차 창(350--500초): EWMA trust가 대략 550 수준까지 하락한다.  
  cost penalty는 적용되지만 아직 parent 후보에서 완전히 제외되지는 않는다.
  \item 2차 창(500--650초): halving decay가 누적되면서 trust가
  $\tau_\text{join}=450$ 이하로 하락한다. 이 시점부터 공격자 노드가 parent 후보 집합에서 제거되기 시작한다.
  \item Recovery 구간(650--900초): 공격자 배제가 충분히 적용되고,
  트래픽이 대안 경로로 재수렴하면서 PDR이 상승한다.
\end{itemize}

이 현상은 전체 30개 seed 중 29개에서 관찰되었으며, 평균 상승폭은 $+5.2$\,pp였다.

반면 SMTrust에서는 이러한 효과가 관찰되지 않는다
(recovery 대비 during 변화량 $-0.05$\,pp).
이는 SMTrust가 공격자 노드를 parent 후보 집합에서 실질적으로 제거하지 못하기 때문이다
(Section~\ref{sec:smtrust}의 TI floor 분석 참조).

\begin{figure}[t]
  \centering
  \includegraphics[width=\linewidth]{fig5_trust_trace.pdf}
  \caption{단일 seed에서의 trust trace. 공격자(attacker) 노드의 $T_\text{fwd}$가
  350초 이후 급격히 하락하여 $\tau_\text{join}=0.45$ 아래로 내려가고,
  이후 recovery 구간에서 대안 경로로의 전환이 이루어지는 과정을 보여준다.}
  \label{fig:trust_trace}
\end{figure}

% ============================================================
\section{V2: 오탐(FPR) 실험}
% ============================================================

공격자가 존재하지 않는 환경에서, TA-BRPL이 정상 노드를 부당하게 격리하는지 확인하기 위해
5-seed 실험을 수행하였다. 이때 기존 공격자 노드는 일반 sender로 대체하였다.

\subsection{PDR 요약과 seed 편차}

\begin{table}[t]
  \centering
  \caption{공격이 없는 환경에서의 PDR (5 seeds)}
  \label{tab:noattack}
  \begin{tabular}{lrrr}
    \toprule
    프로토콜 & Pre & During & Recovery \\
    \midrule
    RPL (no attack)     & 1.0000 & 1.0000 & 1.0000 \\
    BRPL (no attack)    & 1.0000 & 0.9994 & 1.0000 \\
    TA-BRPL (no attack) & 0.9817 & 0.8924 & 0.8076 \\
    \bottomrule
  \end{tabular}
\end{table}

TA-BRPL\_NOATTACK의 during/recovery PDR은 0.8924/0.8076으로 여전히 baseline보다 낮다.
그러나 이는 전형적인 의미의 FPR(정상 노드의 blacklist)이라기보다,
\textbf{합성 혼잡 버스트(200--300초)로 인해 $T_\text{hon}$이 혼잡 노드에 패널티를 부여한 결과}에 가깝다.
동시에 seed 간 편차도 크다. 예를 들어 seed 1은 1.0000/1.0000/1.0000을 유지했지만,
seed 2는 0.9882/0.7048/0.6667까지 하락하였다.
즉 no-attack 저하는 평균값 하나보다도 \emph{특정 seed에서의 불안정성}으로 해석하는 것이 더 적절하다.
이 해석은 V3 분리 실험(Section~\ref{sec:v3}) 및 이후의 seed-level failure analysis와 일관된다.

\begin{table}[t]
  \centering
  \caption{TA-BRPL no-attack seed별 PDR 편차 (5 seeds)}
  \label{tab:noattack-seeds}
  \begin{tabular}{lrrr}
    \toprule
    Seed & Pre & During & Recovery \\
    \midrule
    1 & 1.0000 & 1.0000 & 1.0000 \\
    2 & 0.9882 & 0.7048 & 0.6667 \\
    3 & 0.9735 & 0.9143 & 0.7429 \\
    4 & 0.9882 & 0.8976 & 0.7111 \\
    5 & 0.9588 & 0.9452 & 0.9175 \\
    \bottomrule
  \end{tabular}
\end{table}

\begin{table}[t]
  \centering
  \caption{TA-BRPL no-attack 안정성 진단 요약 (5-seed 평균)}
  \label{tab:noattack-diagnostics}
  \begin{tabular}{lrrrrrr}
    \toprule
    Phase & False BL & Escape & Churn & Allowed cand. & Excl. ratio & $T_\text{ewma}$ \\
    \midrule
    Pre      & 0.4 & 8.4 & 63.8 & 5.91 & 0.0219 & 556.9 \\
    During   & 1.4 & 6.0 & 74.2 & 5.72 & 0.0532 & 578.4 \\
    Recovery & 0.2 & 4.2 & 47.8 & 5.71 & 0.0566 & 594.5 \\
    \bottomrule
  \end{tabular}
\end{table}

\subsection{안정성 진단: blacklist보다 churn}

Table~\ref{tab:noattack-diagnostics}는 no-attack 저하가 단순한 의미의
``대량 false blacklist''라기보다, \textbf{suspect 상태 증가와 parent churn/penalty의 누적 비용}으로
해석되어야 함을 보여준다. 평균 false blacklist count는 during/recovery에서 각각 1.4/0.2에 그쳤고,
trust tier의 black 비율도 사실상 0이었다. 반면 suspect 비율은 pre의 0.816에서 during 0.571,
recovery 0.582 수준으로 넓게 분포했고, parent churn은 pre/during/recovery에서 각각 63.8/74.2/47.8회였다.
즉 PDR 저하의 주원인은 정상 노드의 즉시 격리보다는,
\emph{정상 parent에 대한 반복적인 패널티와 경로 재평가 비용}으로 보는 편이 타당하다.

\begin{figure}[t]
  \centering
  \includegraphics[width=\linewidth]{fig9_noattack_diagnostics.pdf}
  \caption{TA-BRPL no-attack 안정성 진단. blacklist는 많지 않지만, suspect/escape/churn과
  candidate 제외 비율이 during/recovery에서 함께 증가하며 PDR 저하와 동행한다.}
  \label{fig:noattack-diagnostics}
\end{figure}

\begin{figure}[t]
  \centering
  \includegraphics[width=\linewidth]{fig7_noattack_seed_spread.pdf}
  \caption{TA-BRPL no-attack seed별 PDR 편차. 평균 during/recovery PDR은 0.8924/0.8076이지만,
  실제로는 seed 1의 완전 안정(1.0/1.0/1.0)부터 seed 2의 급락(0.7048/0.6667)까지
  불안정성 폭이 크다.}
  \label{fig:noattack-seed-spread}
\end{figure}

\subsection{Route-level 증거}

\begin{figure}[t]
  \centering
  \includegraphics[width=\linewidth]{fig4_route_exposure_timeseries.pdf}
  \caption{공격자 노드를 경유하는 라우트 비율의 시계열(30-seed 평균).
  TA-BRPL은 공격 시작(350초) 이후 가장 빠르게 공격자 노드 노출을 감소시키며,
  recovery 구간에서는 거의 0에 수렴한다.}
  \label{fig:exposure_ts}
\end{figure}

추가적인 route-level 진단도 같은 해석을 뒷받침한다.
30-seed 평균에서 attacker route share는 pre/during/recovery에 대해
0.2081/0.1564/0.0927로 감소하였다.
이때 sinkhole attraction(node 18 share)은 0.0653/0.0203/0.0369,
blackhole aggregate share(node 2/3/4)는 0.1429/0.1360/0.0559였다.
동시에 average hop count는 2.995에서 4.015, 6.241로 증가하고,
top-parent concentration은 0.1752에서 0.1242, 0.1015로 낮아졌다.
즉 TA-BRPL의 recovery 이득은 특정 정상 parent로의 과도한 집중보다는,
\textbf{공격자 경유 비율 감소와 더 분산된 대안 parent 집합 사용}에서 온 것으로 해석할 수 있다.

\begin{figure}[t]
  \centering
  \includegraphics[width=\linewidth]{fig11_route_phase_summary.pdf}
  \caption{Phase별 route-level summary. node 18 및 blackhole parent adoption은 시간이 지나며 감소하고,
  unique parent 수는 회복 구간에서 다시 증가한다.}
  \label{fig:route-phase-summary}
\end{figure}

% ============================================================
\section{V3: 혼잡 vs. 공격 분리 실험}
\label{sec:v3}
% ============================================================

TA-BRPL의 핵심 주장인 ``단순 PDR 하락 $\ne$ 공격''을 검증하기 위해,
혼잡과 공격을 각각 독립적으로 on/off한 4개 시나리오를 구성하여 실험하였다.

\begin{table}[t]
  \centering
  \caption{V3 혼잡--공격 분리 실험 결과 (5 seeds, TA-BRPL)}
  \label{tab:v3}
  \begin{tabular}{lllrrr}
    \toprule
    시나리오 & 혼잡 & 공격 & Pre & During & Recovery \\
    \midrule
    C1: Baseline    & \texttimes & \texttimes & 0.961 & 0.793 & 0.875 \\
    C2: 혼잡만      & \checkmark & \texttimes & 0.980 & 0.857 & 0.789 \\
    C3: 공격만      & \texttimes & \checkmark & 0.994 & 0.379 & 0.595 \\
    C4: 혼잡+공격   & \checkmark & \checkmark & 0.999 & 0.417 & 0.353 \\
    \bottomrule
  \end{tabular}
\end{table}

핵심 비교는 \textbf{C2(0.857) $\gg$ C3(0.379)}이다.
혼잡만 존재하는 경우에는 $T_\text{hon}$이 혼잡 노드에 적절한 패널티를 부여하지만,
전체 PDR은 상대적으로 크게 악화되지 않는다. 반면 공격만 존재하는 경우에는
$T_\text{fwd}$가 급락하면서 조기 blacklist가 유도되고, PDR이 크게 하락한다.
이 대비는 ``TA-BRPL이 혼잡과 공격을 구분한다''는 주장의 실험적 근거를 제공한다.

\begin{figure}[t]
  \centering
  \includegraphics[width=\linewidth]{fig3_attack_tradeoff.pdf}
  \caption{V3 혼잡--공격 분리 실험 결과. C2(혼잡만)와 C3(공격만)의 during PDR 차이가
  뚜렷하게 나타나며, $T_\text{hon}$과 $T_\text{fwd}$가 서로 다른 원인을 탐지함을 보여준다.}
  \label{fig:attack_tradeoff}
\end{figure}

한편, C1(혼잡 없음, 공격 없음)에서도 during PDR이 0.793까지 하락한다.
이는 공격 때문이 아니라, TA-BRPL 자체의 trust penalty 기반 경로 재평가 과정에서
정상 노드에 대해 escape/persistence 패널티가 발동하는 비용으로 해석된다.

% ============================================================
\section{V5: EWMA 파라미터 민감도}
% ============================================================

이 절의 결과는 본선 30-seed 기본 TA-BRPL 설정을 기준으로 수행한
\emph{초기 민감도 실험}이다. 따라서 이후 손실 robustness tuning에서 도입된
PRR 정규화, $J580/W630$, relative filtering 변형과는 직접 비교하지 않는다.

$\lambda_\downarrow$와 $\lambda_\uparrow$의 선택을 정당화하기 위해,
4개의 변형을 각각 5-seed씩 실험하였다(Table~\ref{tab:lambda}).

\begin{table}[t]
  \centering
  \caption{EWMA $\lambda$ 민감도 분석 (5 seeds)}
  \label{tab:lambda}
  \begin{tabular}{llrr}
    \toprule
    변형 & $(\lambda_\downarrow,\,\lambda_\uparrow)$ & During PDR & Recovery PDR \\
    \midrule
    \textbf{기본값 (TA-BRPL)} & $(0.5,\;0.7)$ & \textbf{0.908} & \textbf{0.959} \\
    LAMBDA\_FAST              & $(0.1,\;0.7)$ & 0.384 & 0.444 \\
    LAMBDA\_SLOW              & $(0.4,\;0.7)$ & 0.474 & 0.490 \\
    LAMBDA\_FAST\_RECOVERY    & $(0.2,\;0.5)$ & 0.517 & 0.561 \\
    LAMBDA\_SLOW\_RECOVERY    & $(0.2,\;0.9)$ & 0.534 & 0.685 \\
    \bottomrule
  \end{tabular}
\end{table}

모든 변형은 기본값보다 현저히 낮은 성능을 보였다. 그 해석은 다음과 같다.
\begin{itemize}
  \item \textbf{LAMBDA\_FAST ($\lambda_\downarrow=0.1$)}: 신규 관찰에 90\%의 가중치를 부여하므로
  EWMA 분산이 급격히 증가하고, 정상 노드도 단일 관찰 창에서 FP blacklist될 수 있다.
  그 결과 during PDR이 0.384까지 하락하여 가장 나쁜 성능을 보였다.
  \item \textbf{LAMBDA\_SLOW ($\lambda_\downarrow=0.4$)}: 이름과 달리 기본값(0.5)보다
  신규 관찰 비중이 더 크므로 오히려 더 반응적이다(60\% vs.\ 50\%).
  따라서 FP 비율이 상승한다.
  \item \textbf{FAST/SLOW\_RECOVERY}: 두 변형 모두 $\lambda_\downarrow=0.2$를 사용하므로
  이미 과반응 영역에 진입한 상태다. 따라서 회복 속도의 차이보다 공격 중 FP 증가의 영향이 더 크게 나타난다.
\end{itemize}

결론적으로, $\lambda_\downarrow=0.5$는 공격 반응성과 EWMA 안정성 사이의 균형점으로 해석된다.
특히 $\lambda_\downarrow < 0.5$인 경우에는 분산이 과도하게 증가하여 성능 저하를 초래한다.

% ============================================================
\section{V6: 임계값 민감도}
% ============================================================

이 절 역시 본선 기본 설정의 절대 threshold 구조를 기준으로 수행한
\emph{초기 민감도 실험}이다. Section~\ref{tab:loss} 이후의 loss-ladder tuning은
별도의 robustness 확장으로 해석해야 한다.

$(\tau_\text{warn}, \tau_\text{join}, \tau_\text{black})$의 선택을 정당화하기 위해,
3개 변형을 각각 5-seed씩 실험하였다(Table~\ref{tab:thresh}).

\begin{table}[t]
  \centering
  \caption{임계값 민감도 분석 (5 seeds). 기본값: $(0.70,\;0.45,\;0.25)$}
  \label{tab:thresh}
  \begin{tabular}{lp{2.2cm}rr}
    \toprule
    변형 & $(\tau_w,\;\tau_j,\;\tau_b)$ & Pre PDR & During PDR \\
    \midrule
    \textbf{기본값} & $(0.70,\;0.45,\;0.25)$ & 0.999 & \textbf{0.908} \\
    THRESH\_STRICT  & $(0.80,\;0.60,\;0.40)$ & 0.914 & 0.586 \\
    THRESH\_RELAXED & $(0.70,\;0.50,\;0.30)$ & 0.993 & 0.494 \\
    THRESH\_JOINLOW & $(0.75,\;0.40,\;0.25)$ & 0.999 & 0.750 \\
    \bottomrule
  \end{tabular}
\end{table}

\begin{itemize}
  \item \textbf{THRESH\_STRICT}: $\tau_\text{join}=0.60$이 지나치게 높아,
  네트워크 수렴 과정에서 자연스럽게 forwarding rate가 낮은 정상 노드까지
  후보에서 제외된다. 그 결과 pre PDR이 이미 0.914까지 하락한다.
  \item \textbf{THRESH\_RELAXED}: $\tau_\text{join}=0.50$으로 상승하면
  trust가 $[0.45,0.50)$ 구간에 있는 노드가 더 오래 후보로 남게 되어 탐지가 지연된다.
  결과적으로 during PDR이 0.494로 크게 악화된다.
  \item \textbf{THRESH\_JOINLOW}: $\tau_\text{join}=0.40$으로 낮추면
  suspect 구간이 넓어지고(0.40$\sim$0.75), 배제 전 패널티 기간이 길어진다.
  다만 완전 배제까지의 시간이 증가하여 during PDR은 0.750 수준에 머문다.
\end{itemize}

결론적으로 기본값 $(0.70, 0.45, 0.25)$는 EWMA 파라미터와 함께 설계된
실질적 최적점으로 해석된다. 특히 $\tau_\text{warn} - \tau_\text{join} \ge 0.20$의
간격을 확보하는 것이 중요하다.

% ============================================================
\section{절제 실험: 3성분 필요성 검증}
% ============================================================

각 trust 성분의 기여를 독립적으로 확인하기 위해,
단일 성분 또는 2성분 변형을 각각 5-seed씩 실험하였다(Table~\ref{tab:ablation}).

\begin{table}[t]
  \centering
  \caption{성분 절제 실험 결과 (5 seeds)}
  \label{tab:ablation}
  \begin{tabular}{lllrrr}
    \toprule
    변형 & 성분 & Pre & During & Recovery \\
    \midrule
    TA-BRPL (전체) & $T_f + T_c + T_h$ & 0.999 & \textbf{0.908} & \textbf{0.959} \\
    FWD+CTRL       & $T_f + T_c$        & 0.926 & 0.660 & 0.516 \\
    FWD only       & $T_f$ only         & 0.898 & 0.719 & 0.731 \\
    \midrule
    RPL (기준)     & 없음               & 1.000 & 0.876 & 0.875 \\
    \bottomrule
  \end{tabular}
\end{table}

흥미롭게도, \textbf{두 변형 모두 RPL보다 낮은 성능을 보였다}.
예를 들어 FWD only의 during PDR(0.719)은 RPL의 during PDR(0.876)보다 낮다.

그 이유는 다음과 같다.
\begin{enumerate}
  \item FWD only에서는 pre PDR이 이미 0.898로 하락한다.  
  이는 네트워크 수렴 초기의 자연스러운 forwarding rate 저하가 FP blacklist로 이어지기 때문이다.
  \item FWD+CTRL에서는 pre PDR이 0.926으로 다소 개선되지만,
  DIO trickle 정착 구간에서 DIO 빈도 이상 역시 FP를 추가로 유발한다.
  \item $T_\text{hon}$이 있어야 이러한 FP가 억제된다.  
  혼잡 구간에서 정상 노드의 $T_\text{hon} \approx 1.0$이 EWMA를 안정화하여,
  pre-attack FP를 줄이는 역할을 한다.
\end{enumerate}

즉, 세 trust 성분은 선택적으로 더하는 부가 요소가 아니라,
\textbf{함께 사용될 때에만 의미를 갖는 상보적 구성}으로 해석할 수 있다.
특히 $T_\text{hon}$이 없으면 $T_\text{fwd}$ 단독 사용은 오히려 RPL보다 더 나쁜 결과를 낳는다.

\begin{figure}[t]
  \centering
  \includegraphics[width=\linewidth]{fig6_churn_hotspots.pdf}
  \caption{Parent churn 핫스팟 분포. FWD-only 변형에서는 공격자 인근 노드의
  parent 교체 빈도가 급증하며, $T_\text{hon}$을 추가하면 이 현상이 억제된다.}
  \label{fig:churn}
\end{figure}

% ============================================================
\section{손실률 로버스트니스와 튜닝 진행}
% ============================================================

v0.3.2까지의 핵심 수정은 $T_\text{fwd}$ 분모를 단순 $S_{ij}$가 아니라
ETX 기반 기대 forwarding 수 $E_{ij}$로 바꾸는 것이었다.
이 변경은 저손실 구간의 FP를 완화하는 데는 효과적이었지만,
중간 이상의 손실률에서는 trust 제어가 여전히 공격자와 정상 이웃을 충분히 분리하지 못했다.
이에 v0.3.3에서는 5-seed loss ladder(10\%, 30\%, 50\%)를 대상으로
다음의 tuning progression을 순차적으로 수행하였다.
\begin{enumerate}
  \item ETX 기반 PRR 정규화(v0.3.2)
  \item 절대 threshold 재튜닝: $(\tau_\text{join}, \tau_\text{warn})=(580,630)$
  \item soft-penalty / escape 조건 sweep
  \item median 기반 relative trust filtering
\end{enumerate}

이 절은 이 순서를 그대로 따라간다. 먼저 absolute threshold와 soft penalty가 왜 실패했는지를 보이고,
그 다음 attacker/benign trust 분리도와 seed dependence를 통해 기준 모델 내부의 실패 패턴을 정리한 뒤,
마지막으로 relative filtering이 어느 부분을 완화했는지와 sliding-window 대안이 왜 제외되었는지를 설명한다.

\begin{table}[t]
  \centering
  \caption{Loss ladder에서의 대표 비교 결과 (5 seeds)}
  \label{tab:loss}
  \begin{tabular}{llrr}
    \toprule
    손실률 & 프로토콜/변형 & During PDR & Recovery PDR \\
    \midrule
    \multirow{4}{*}{10\%}
      & BRPL                         & \textbf{0.838} & \textbf{0.806} \\
      & TA-BRPL (기본)              & 0.778 & 0.767 \\
      & TA-BRPL + J580/W630         & 0.643 & 0.627 \\
      & TA-BRPL + REL\_020          & 0.637 & 0.664 \\
    \midrule
    \multirow{4}{*}{30\%}
      & BRPL                         & \textbf{0.694} & \textbf{0.683} \\
      & TA-BRPL (기본)              & 0.610 & 0.547 \\
      & TA-BRPL + J580/W630         & 0.609 & 0.644 \\
      & TA-BRPL + REL\_010          & 0.587 & 0.658 \\
    \midrule
    \multirow{4}{*}{50\%}
      & BRPL                         & \textbf{0.588} & \textbf{0.567} \\
      & TA-BRPL (기본)              & 0.508 & 0.505 \\
      & TA-BRPL + J580/W630         & 0.503 & 0.495 \\
      & TA-BRPL + REL\_050          & 0.539 & 0.533 \\
    \bottomrule
  \end{tabular}
\end{table}

Table~\ref{tab:loss}는 세 가지 중요한 사실을 보여준다.
첫째, ETX 기반 PRR 정규화만으로는 10\% 손실에서는 개선이 있었지만,
30\%/50\% 손실에서 BRPL과의 격차를 줄이기에 충분하지 않았다.
둘째, 절대 threshold를 $(580,630)$으로 올리면 attacker share는 줄지만
usable parent 수가 감소하여 10\% 손실 구간 PDR가 오히려 하락할 수 있었다.
셋째, relative filtering은 이 절대 threshold 구조보다
``attacker share 감소 + usable parent 유지''의 균형을 더 잘 달성하였다.

\subsection{절대 threshold 재튜닝의 진단}

threshold sweep 결과, $J580/W630$은 손실 ladder 전체에서 가장 균형적인 절충점이었다.
예를 들어 30\% 손실에서는 during/recovery PDR이 0.609/0.644였고,
attacker share는 0.208/0.207로 기본 TA-BRPL의 0.264/0.277보다 낮아졌다.
그러나 10\% 손실에서는 during/recovery PDR이 0.643/0.627로 떨어져,
BRPL(0.838/0.806)과의 격차가 오히려 커졌다.

soft penalty를 더 강하게 주면 attacker share는 추가로 감소했지만,
정상 parent까지 함께 pruning되는 문제가 드러났다.
대표적으로 10\% 손실에서 \texttt{SFT\_BASE}는 attacker share 0.172/0.171,
usable parents 2.69/2.70, PDR 0.643/0.627이었고,
\texttt{SFT\_G2}는 attacker share를 0.165/0.173으로 더 낮췄지만
usable parents가 2.54/2.52로 줄어들며 PDR이 0.574/0.629로 악화되었다.
이 결과는 ``탐지 신호는 둔한데 제어 로직은 민감하다''는 구조적 불균형을 보여준다.

\subsection{원인 분석: A $>$ C $\gg$ B}

10\% 손실 구간에서 ``attacker share는 낮아졌는데 PDR이 더 낮은'' 현상을 설명하기 위해,
세 가지 가설을 비교하였다.
\begin{itemize}
  \item \textbf{A. 정상 parent의 FP성 pruning}: 공격자뿐 아니라 정상 이웃도 trust 하향에 함께 노출되어
  후보 집합이 과도하게 줄어드는 경우
  \item \textbf{B. backpressure 교란}: trust penalty가 queue-aware ranking을 왜곡하여
  특정 parent로 트래픽이 과집중되는 경우
  \item \textbf{C. 경로 수렴 실패}: 공격자 배제 후 대안 경로 재수렴이 길어지는 경우
\end{itemize}

현재 로그는 A가 주원인이고, C가 보조 원인이며, B의 증거는 약함을 보여준다.
정상 이웃의 $T_\text{fwd}$와 $T_\text{ewma}$도 공격자와 함께 하향했고,
usable parent 수가 줄어드는 반면 top-parent concentration은 뚜렷하게 증가하지 않았다.
즉, soft penalty 강화의 부작용은 ``공격자를 더 강하게 누르기보다
정상 후보까지 같이 줄이는 것''으로 해석하는 것이 타당하다.

\subsection{기준 모델 내부 진단: Seed 3 vs.\ Seed 28}

기준 모델 자체의 불안정성을 더 구체적으로 보기 위해,
30-seed baseline 중 가장 불안정한 \texttt{seed 3}과 가장 안정적인 \texttt{seed 28}을 직접 비교하였다.
두 seed는 동일한 알고리즘과 파라미터를 사용하지만 결과는 크게 달랐다.
\texttt{seed 3}의 PDR은 pre/during/recovery = 0.9877/0.2430/0.2661이었고,
\texttt{seed 28}은 1.0000/0.5699/0.7688이었다.

\begin{figure}[t]
  \centering
  \fbox{%
  \begin{minipage}{0.94\linewidth}
  \centering
  \small
  \begin{tabular}{lrr}
    \toprule
    지표 & seed 3 (불안정) & seed 28 (안정) \\
    \midrule
    During PDR & 0.2430 & 0.5699 \\
    Recovery PDR & 0.2661 & 0.7688 \\
    During allowed candidates (mean) & 2.89 & 4.13 \\
    Recovery allowed candidates (mean) & 1.58 & 3.23 \\
    Recovery $attacker$ route share & 0.182 & 0.037 \\
    Recovery attacker 18 $T_\text{ewma}$ (mean) & 649.3 & 494.5 \\
    During blacklist count & 2 & 0 \\
    Recovery blacklist count & 2 & 0 \\
    \bottomrule
  \end{tabular}
  \end{minipage}}
  \caption{기준 TA-BRPL 내부의 seed-dependent failure pattern. 불안정한 seed 3에서는
  attacker trust가 충분히 내려가지 않는 동시에 benign parent pruning이 겹치며,
  recovery 단계에서 candidate set이 급격히 붕괴한다.}
  \label{fig:seed3-vs-28}
\end{figure}

Figure~\ref{fig:seed3-vs-28}은 중요한 두 가지 사실을 보여준다.
첫째, 흔들림의 핵심은 escape 이벤트 수 자체가 아니다.
\texttt{seed 28}은 during-attack 구간에서 escape가 더 많이 발생했음에도 recovery가 잘 유지되었다.
반면 \texttt{seed 3}은 escape 수가 상대적으로 적었지만 recovery가 크게 무너졌다.
즉 병목은 ``escape의 빈도''보다 \textbf{escape 이후에도 유지되는 parent 후보 집합의 질}에 있다.

둘째, 최악의 경우는 \textbf{attacker trust가 충분히 빨리 내려가지 않는 동시에 benign pruning이 겹치는 조합}이다.
\texttt{seed 3}에서는 attacker 18의 $T_\text{ewma}$가 recovery에서도 평균 649.3으로 높게 남았고,
동시에 allowed candidate 수는 1.58까지 감소하였다.
그 결과 공격자 완전 배제가 일어나지 않은 상태에서 정상 parent set만 먼저 줄어들며,
PDR recovery가 0.2661까지 하락하였다.
반대로 \texttt{seed 28}에서는 recovery 단계 attacker route share가 0.037까지 줄고
candidate 수가 3.23으로 유지되어 recovery PDR 0.7688을 달성하였다.

\subsection{Relative Trust Filtering의 효과}

이 진단을 바탕으로 v0.3.3에서는 절대 임계값 대신
이웃 trust 분포의 median을 기준으로 후보를 자르는 relative filtering을 구현하였다.
핵심 목적은 손실 환경에서 전체 trust 수준이 함께 낮아질 때,
정상 노드들이 median 근방에서 살아남도록 하는 것이다.

5-seed sweep에서 손실률별 best margin은 서로 달랐다.
\begin{itemize}
  \item 10\% 손실: \texttt{REL\_020}, PDR 0.637/0.664, attacker share 0.206/0.184
  \item 30\% 손실: \texttt{REL\_010}, PDR 0.587/0.658, attacker share 0.211/0.183
  \item 50\% 손실: \texttt{REL\_050}, PDR 0.539/0.533, attacker share 0.219/0.213
\end{itemize}

특히 \texttt{REL\_050}은 세 손실 구간 모두에서 비교적 안정적인 균형점을 제공했다.
10\%/30\%/50\% 손실에서 PDR은 각각 0.656/0.639, 0.614/0.639, 0.539/0.533이었고,
attacker share는 0.216/0.208, 0.217/0.188, 0.219/0.213으로 유지되었다.
즉, 단일 margin으로 모든 구간을 최적화하지는 못했지만,
absolute threshold 대비 ``parent 후보 유지'' 측면에서 더 일관된 거동을 보였다.

\subsection{Trust 분리도와 threshold crossing 진단}

기준 모델 내부의 흔들림이 정말 ``attacker under-suppression + benign overlap''인지 보기 위해,
attacker/benign 이웃의 trust 성분 분포를 직접 비교하였다.
Table~\ref{tab:trust-separability}의 overlap coefficient는 두 분포가 얼마나 겹치는지를 나타내며,
1에 가까울수록 분리력이 약함을 의미한다.

\begin{table}[t]
  \centering
  \caption{Attacker vs.\ benign trust separability (30-seed 평균)}
  \label{tab:trust-separability}
  \begin{tabular}{llrrr}
    \toprule
    Phase & Metric & Attacker mean & Benign mean & Overlap \\
    \midrule
    During   & $T_\text{fwd}$  & 345.7 & 409.1 & 0.748 \\
    During   & $T_\text{ctrl}$ & 978.2 & 951.1 & 0.816 \\
    During   & $T_\text{hon}$  & 1000.0 & 1000.0 & 1.000 \\
    During   & $T_\text{ewma}$ & 462.5 & 517.4 & 0.739 \\
    \midrule
    Recovery & $T_\text{fwd}$  & 329.9 & 378.4 & 0.725 \\
    Recovery & $T_\text{ctrl}$ & 909.9 & 792.3 & 0.385 \\
    Recovery & $T_\text{hon}$  & 1000.0 & 1000.0 & 1.000 \\
    Recovery & $T_\text{ewma}$ & 438.9 & 472.3 & 0.500 \\
    \bottomrule
  \end{tabular}
\end{table}

표가 보여주는 핵심은 세 가지다.
첫째, blackhole 탐지에 가장 직접적으로 연결된 $T_\text{fwd}$조차 during/recovery overlap이
0.748/0.725로 여전히 크다. 둘째, 현재 실험에서는 $T_\text{hon}$이 attacker와 benign 모두에서
거의 1.0으로 포화되어, no-attack 혼잡 완화에는 도움을 주더라도 공격자 분리 신호로는 거의 작동하지 않는다.
셋째, aggregate trust인 $T_\text{ewma}$도 overlap 0.739/0.500 수준을 보여,
일부 seed에서는 공격자와 정상 노드가 여전히 같은 threshold band에 겹쳐질 수 있음을 시사한다.

\begin{figure}[t]
  \centering
  \includegraphics[width=\linewidth]{fig10_trust_separability.pdf}
  \caption{Attacker vs.\ benign trust separability. $T_\text{fwd}$와 $T_\text{ewma}$는
  attacker 쪽이 더 낮지만, overlap이 여전히 커서 일부 seed에서는 threshold band가 겹친다.}
  \label{fig:trust-separability-plot}
\end{figure}

threshold crossing 시간도 이 해석과 일치한다.
attacker 이웃은 attack 시작 직후인 390001\,ms 부근에서
$\tau_\text{warn}$과 $\tau_\text{join}$을 통과하는 경우가 많았지만,
$\tau_\text{black}$ 아래로 내려간 사례는 관측되지 않았다.
반면 benign 이웃도 $\tau_\text{join}$을 median 510001\,ms 수준에서 자주 통과했고,
$\tau_\text{black}$ 아래로 내려간 사례도 소수 존재하였다.
즉 현재 기준 모델의 실패는 ``공격자만 충분히 빨리 낮추지 못함''과
``정상 이웃도 threshold 아래로 일부 함께 끌려감''이 동시에 존재하는 문제로 볼 수 있다.

\begin{figure}[t]
  \centering
  \includegraphics[width=0.78\linewidth]{fig12_threshold_crossing.pdf}
  \caption{Median threshold crossing time. attacker는 $\tau_\text{warn}$ 아래로는 빠르게 내려가지만,
  $\tau_\text{black}$까지는 거의 도달하지 못하고, benign 역시 $\tau_\text{join}$ 아래로 일부 함께 내려간다.}
  \label{fig:threshold-crossing}
\end{figure}

\subsection{Population-level Seed Dependence}

seed 3 vs.\ 28은 개별 사례였지만, 전체 30개 seed 분포를 보면 같은 패턴이 population level에서도 나타난다.
During/recovery PDR는 각각 평균 0.530/0.405였고, 상위 5개 seed와 하위 5개 seed를 비교하면
recovery PDR는 0.616 vs.\ 0.219로 큰 차이를 보였다.
이 차이는 route concentration보다는 candidate set과 attacker trust 쪽 지표와 더 강하게 연결되었다.

\begin{figure}[t]
  \centering
  \includegraphics[width=\linewidth]{fig8_seed_failure_correlation.pdf}
  \caption{30-seed population-level failure pattern. Recovery allowed candidate 수가 많을수록 recovery PDR가 높고,
  attacker 18의 recovery $T_\text{ewma}$가 높을수록 recovery PDR가 낮아지는 경향이 나타난다.}
  \label{fig:seed-failure-corr}
\end{figure}

\begin{table}[t]
  \centering
  \caption{Worst 5 vs.\ Best 5 seeds 비교 (기준 TA-BRPL)}
  \label{tab:worstbest5}
  \begin{tabular}{lrr}
    \toprule
    지표 & Worst 5 mean & Best 5 mean \\
    \midrule
    During PDR & 0.409 & 0.584 \\
    Recovery PDR & 0.219 & 0.616 \\
    During allowed candidates & 3.79 & 4.26 \\
    Recovery allowed candidates & 2.53 & 3.20 \\
    During attacker share & 0.171 & 0.151 \\
    Recovery attacker share & 0.153 & 0.068 \\
    Recovery attacker 18 $T_\text{ewma}$ & 543.1 & 444.0 \\
    During blacklists & 0.2 & 0.0 \\
    Recovery blacklists & 0.6 & 0.0 \\
    During escapes & 16.6 & 26.6 \\
    \bottomrule
  \end{tabular}
\end{table}

상관 분석도 같은 결론을 지지한다.
During PDR는 during/recovery allowed candidate 수와 각각 $r=0.830$, $r=0.825$의 강한 양의 상관을 보였고,
recovery attacker 18 $T_\text{ewma}$와는 $r=-0.772$의 강한 음의 상관을 보였다.
반면 recovery top-parent share와 during PDR의 상관은 $r=0.154$로 약했다.
즉 흔들린 seed의 공통점은 ``어느 한 정상 parent로의 과집중''이 아니라,
\textbf{candidate collapse와 attacker under-suppression의 동시 발생}으로 보는 것이 타당하다.

% ============================================================
\section{전체 파라미터 목록}
% ============================================================

\begin{table}[t]
  \centering
  \caption{TA-BRPL 주요 파라미터 (\texttt{project-conf.h} 현재 기본값 기준)}
  \label{tab:params}
  \begin{tabular}{lll}
    \toprule
    파라미터 & 심벌 & 기본값 \\
    \midrule
    \multicolumn{3}{l}{\textit{임계값}} \\
    Warn 임계값       & $\tau_\text{warn}$    & 700 \\
    Join 임계값       & $\tau_\text{join}$    & 450 \\
    Blacklist 임계값  & $\tau_\text{black}$   & 250 \\
    초기 신뢰         & $T_\text{init}$       & 500 \\
    \midrule
    \multicolumn{3}{l}{\textit{EWMA}} \\
    감소 람다         & $\lambda_\downarrow$  & 0.2 \\
    정상 람다         & $\lambda_\uparrow$    & 0.7 \\
    업데이트 주기     & $\Delta t$            & 60s \\
    \midrule
    \multicolumn{3}{l}{\textit{PRR 정규화}} \\
    최소 PRR          & $\mathrm{PRR}_{\min}$ & 100/1000 \\
    PRR blend weight  & --                    & 1000 \\
    최대 PRR          & $\mathrm{PRR}_{\max}$ & 1000/1000 \\
    ETX fallback PRR  & --                    & 1000/1000 \\
    \midrule
    \multicolumn{3}{l}{\textit{Relative filtering (기본 비활성)}} \\
    Relative filter enable  & -- & 0 \\
    Relative penalty enable & -- & 0 \\
    Relative margin         & $\Delta_\text{rel}$ & 25 \\
    Relative penalty scale  & $S_\text{rel}$ & 1000 \\
    Relative penalty max    & $P^\text{rel}_\text{max}$ & 400 \\
    \midrule
    \multicolumn{3}{l}{\textit{가중치}} \\
    $T_\text{fwd}$ 가중치   & $w_f$      & 7 \\
    $T_\text{ctrl}$ 가중치  & $w_c$      & 2 \\
    $T_\text{hon}$ 가중치   & $w_h$      & 1 \\
    Laplace $\alpha$        & $\alpha$   & 1 \\
    Laplace $\beta$         & $\beta$    & 1 \\
    Rank 이상 가중치        & $w_1$      & 5 \\
    DIO 이상 가중치         & $w_2$      & 3 \\
    Version 이상 가중치     & $w_3$      & 2 \\
    \midrule
    \multicolumn{3}{l}{\textit{Blacklist / 회복}} \\
    격리 기간               & $t_\text{bl}$     & 120s \\
    해제 후 복원 신뢰       & $T_\text{rst}$    & 450 \\
    해제 cooldown           & $t_\text{cool}$   & 120s \\
    해제 초기 패널티 배율   & $P_\text{rel}$    & $1.6\times$ \\
    지속 패널티 창          & $t_\text{per}$    & 120s \\
    지속 패널티 step        & $\Delta P$        & 250 \\
    지속 패널티 최대값      & $P_\text{max}$    & $2.6\times$ \\
    Escape 발동 시간        & $t_\text{esc}$    & 180s \\
    Escape 패널티 배율      & $P_\text{esc}$    & $3.2\times$ \\
    Escape 신뢰 임계값      & $\tau_\text{esc}$ & 700 \\
    Escape 연속 조건        & --                & 1 update \\
    Escape cooldown         & --                & 0s \\
    Better-parent only      & --                & off \\
    Escape $T_\text{fwd}$ 조건 & --             & $< 450$ \\
    Escape $T_\text{hon}$ 조건 & --             & $> 850$ \\
    Join 최소 지속 시간     & --                & 120s \\
    Black 최소 지속 시간    & --                & 120s \\
    \midrule
    \multicolumn{3}{l}{\textit{실험 환경}} \\
    토폴로지     & -- & GRID $6\times6$, 36 노드 \\
    공격자       & -- & Blackhole $\times3$ + Sinkhole $\times1$ \\
    Seed 수      & -- & 30 \\
    송신 주기    & -- & 30s \\
    시뮬레이션   & -- & 900s \\
    혼잡 유도    & -- & on, 200--300s, 15s interval \\
    \bottomrule
  \end{tabular}
\end{table}

손실 robustness tuning에서 사용한 대표 변형은 기본값을 그대로 바꾸지 않고
\texttt{Makefile} 오버라이드로 실험하였다. 대표적으로 절대 threshold 재튜닝은
$J580/W630$, relative filtering은 \texttt{REL\_010}--\texttt{REL\_100},
soft penalty는 \texttt{SFT\_BASE}, \texttt{SFT\_G2}, \texttt{SFT\_L600} 등으로 분기하였다.

% ============================================================
\section{구현상 한계}
% ============================================================

\textbf{L1 정적 토폴로지}: 모든 실험은 고정 GRID 토폴로지에서 수행되었다.
이동성이 있는 환경에서는 $T_\text{fwd}$의 overhearing 범위 이탈로 인해
FP가 증가할 가능성이 있다.

\textbf{L2 중간 이상 링크 손실}: v0.3.3에서는 10\%/30\%/50\% loss ladder와
절대 threshold, soft penalty, relative filtering tuning까지 5-seed로 수행하였다.
그러나 아직 30-seed robustness 재검증은 완료되지 않았고,
단일 설정이 모든 손실 구간에서 BRPL을 일관되게 상회하지도 못했다.
즉, 현재 데이터는 ``어떤 tuning이 어디서 실패하는가''를 진단하기에는 충분하지만,
``TA-BRPL이 고손실에서도 전반적으로 우월하다''는 강한 주장을 하기에는 아직 부족하다.
또한 sliding-window $T_\text{fwd}$는 무손실 baseline 자체를 크게 훼손하여
본선 대안으로 채택하지 않았다.
기준 모델 내부 분석(Figure~\ref{fig:seed3-vs-28})은
성능 저하가 단순한 escape 과민반응이 아니라,
``attacker trust 미하강 + benign pruning''이 동시에 발생하는 seed-dependent 최악 조합에서
발생함을 보여준다.

\textbf{L3 제한된 공격 유형}: 현재 평가는 Blackhole + rank sinkhole에 한정되어 있다.
Colluding sinkhole, version attack, Sybil attack 등은 아직 평가하지 않았다.

\textbf{L4 도청 불완전성}: IP 레이어 도청 훅은 MAC 필터링으로 폐기된 패킷을 감지하지 못한다.
따라서 손실 환경에서 $F_{ij}$ 관측치가 구조적으로 과소계산될 가능성이 있다.
현재 UDGM에서는 이 문제가 부분적으로 가려져 있지만, 실제 하드웨어나 더 현실적인 radio 모델에서는
$T_\text{fwd}$ 분리력이 더 달라질 수 있다.

\textbf{L5 SMTrust DIO 확장}: TMRT 성분은 비표준 ICMPv6 옵션(0xFE, 0x02)을 사용한다.
따라서 RFC 호환성이 없다.

\textbf{L6 부동소수점}: 현재 Cooja 구현에서는 기하평균 계산에 glibc의
\texttt{pow()}를 사용한다. 실제 MCU(Cortex-M0+) 배포를 고려한다면,
정수 기반 lookup table 방식으로 대체할 필요가 있다.

\textbf{L7 정적 임계값의 한계}: 절대 threshold와 단일 relative margin 모두
손실률 전 구간에서 공통 최적점을 제공하지 못했다.
현재 결과는 candidate-aware threshold, loss-aware adaptive threshold,
혹은 relative spread(MAD/IQR) 기반 판정이 다음 단계의 유력한 개선 방향임을 시사한다.
동시에 보다 직접적인 남은 문제는,
원래 TA-BRPL 기준 모델 내부에서 왜 특정 seed와 조건에서 attacker $T_\text{fwd}$가 충분히 빨리 내려가지 않는가를
추적하는 것이다. 즉 다음 단계의 진단 초점은 새 모델 추가가 아니라,
기존 echo-credit 기반 관측과 trust-to-control coupling의 상호작용을 해석하는 데 있다.

\textbf{L8 정량 overhead 계측 부족}: 현재 로그는 parent churn, escape 횟수,
candidate 수, route share와 같은 \emph{간접 비용 지표}는 충분히 제공하지만,
DIO/DAO/DIS 개수, byte-level control overhead, queue occupancy, CPU time, RAM/ROM 증가량을
직접 기록하지는 않는다. 따라서 본 문서는 현재 단계에서
``보안/복원력 이득에 어떤 경로 제어 비용이 수반되는가''까지는 분석할 수 있지만,
``정확히 몇 바이트/몇 사이클의 오버헤드인가''까지는 아직 정량화하지 못한다.
향후 버전에서는 control-plane 카운터와 메모리 footprint 로깅을 추가하여
성능 향상이 overhead를 정당화하는지까지 함께 평가할 필요가 있다.

\end{document}
